\begin{ExerciseBlock}
\ExerciseHeading{02}{Mittlere Geschwindigkeit und Momentangeschwindigkeit}{linear_kinematics_sports_car_acceleration}
Die Geschwindigkeit eines Sportwagens, der mit einer Beschleunigung $a = 4 \, m/s^2$ anfährt, soll $50$m hinter dem Startpunkt gemessen werden. Dazu wird der Durchgang des Wagens bei $45$m und $55$m gestoppt und die mittlere Geschwindigkeit $\bar{v}$ bestimmt. Um wieviel Prozent weicht diese mittlere Geschwindigkeit $\bar{v}$ von der tatsächlichen Geschwindigkeit (Momentangeschwindigkeit) bei $50$m ab?
\end{ExerciseBlock}
\UseExerciseSolution{linear_kinematics_sports_car_acceleration}
